\newcommand{\legI}{\textbf{Fig. 1 | O corte de geração fotovoltaica no SIN é sistêmico, de razão energética e incide sobre toda a janela solar.} \textbf{a}, Perfil diurno médio da potência fotovoltaica das usinas com apuração de \textit{constrained-off} no Sistema Interligado Nacional. A área azul é a geração verificada e a área vermelha é a energia cortada, definida como \texttt{max(referência $-$ verificada, 0)} nos intervalos com restrição apurada; a soma define a potência que estaria disponível. Médias sobre todos os intervalos de 30 min do período; horário de Brasília, não corrigido para hora solar local. \textbf{b}, Taxa mensal de corte por subsistema, definida como energia cortada dividida pela energia disponível (verificada + cortada). \textbf{c}, Energia total cortada por razão de restrição, decomposta em origem sistêmica e local; rótulos indicam o total por razão. \textbf{d}, Distribuição espacial das usinas e conjuntos fotovoltaicos com apuração de \textit{constrained-off}; a cor codifica a taxa de corte acumulada e o tamanho, a capacidade instalada. A escala de cor satura no percentil 97 para preservar contraste. Contornos estaduais: IBGE. Dados: ONS, dados abertos, restrição por \textit{constrained-off} de usinas fotovoltaicas e fator de capacidade. Dados-fonte em \texttt{data/processed/fig1\_source\_*.csv}.}
\newcommand{\legII}{\textbf{Fig. 2 | A MMGD brasileira é majoritariamente sub-pixel no Sentinel-2, e a série de conexões carrega uma descontinuidade regulatória.} \textbf{a}, Distribuição acumulada da área do arranjo das instalações fotovoltaicas de micro e minigeração distribuída registradas na ANEEL. Linhas tracejadas marcam a área de um pixel de cada sensor; o rótulo indica a fração de instalações menor que um pixel. \textbf{b}, Fração de instalações que atinge o piso operacional de 4 pixels, por classe de consumo e sensor. \textbf{c}, Conexões mensais; a linha tracejada marca o fim da janela de transição da Lei 14.300/2022 e o marcador, o mês de pico. \textbf{d}, Participação da classe rural no total de instalações de cada unidade da federação. Contornos: IBGE. Dados: ANEEL, relação de empreendimentos de geração distribuída e informações técnicas fotovoltaicas. Dados-fonte em \texttt{data/processed/fig2\_source\_*.csv}.}
\newcommand{\legIII}{\textbf{Fig. 3 | Dois ativos solares em Minas Gerais diferem 3,0$\times$ na taxa de corte, e a diferença não se reproduz como padrão no parque.} \textbf{a}, Perfil diurno médio de potência, em pu da capacidade instalada, na janela comum aos dois ativos. Azul, geração verificada; vermelho, energia cortada. \textbf{b}, Custo Marginal de Operação do subsistema Sudeste ponderado pela energia, durante os intervalos de corte e durante a geração. \textbf{c}, Taxa de corte mensal dos dois ativos. \textbf{d}, Pixels Sentinel-2 ocupados pela área de módulos, comparados à instalação mediana de MMGD; escala logarítmica, área de módulos estimada a 5,12 m\textsuperscript{2}/kWp. \textbf{e}, Taxa de corte contra capacidade renovável variável num raio de 100 km, para os conjuntos com apuração completa na janela de referência da frota; os dois alvos em destaque. Dados: ONS (constrained-off, fator de capacidade, CMO, relacionamento conjunto-usina) e ANEEL (SIGA). Dados-fonte em \texttt{data/processed/fig3\_source\_*.csv}.}
\newcommand{\legIV}{\textbf{Fig. 4 | A queda do fator de capacidade de Sol do Cerrado acompanha a energização do entorno; o controle não mostra o mesmo.} \textbf{a}, Jaíba (MG). Área cinza, capacidade fotovoltaica acumulada em operação num raio de 100 km (eixo esquerdo); linha, fator de capacidade mensal verificado de Sol do Cerrado (eixo direito); faixa vermelha, parcela cortada apurada pelo ONS, disponível apenas a partir de abril de 2024 (linha tracejada). \textbf{b}, São Gonçalo do Abaeté (MG), com a mesma construção para Três Marias 3. \textbf{c}, Fator de capacidade mensal de Sol do Cerrado em 2023 e 2025, mês contra mês, para separar restrição de sazonalidade. \textbf{d}, Potência acumulada de micro e minigeração distribuída nos dois municípios; os 20,4 MW de Jaíba equivalem a 3,0\% da capacidade de Sol do Cerrado. Dados: ONS (fator de capacidade, constrained-off) e ANEEL (SIGA, MMGD). Dados-fonte em \texttt{data/processed/fig4\_source\_*.csv}.}
\newcommand{\legV}{\textbf{Fig. 5 | Normalizada pela irradiância medida sobre a usina, a perda de Sol do Cerrado permanece; o recurso solar não a explica.} \textbf{a}, Razão de desempenho mensal, definida como energia gerada por kW instalado dividida pela irradiância global horizontal diária; tracejado horizontal marca a média de cada ano. \textbf{b}, Irradiância global horizontal mensal sobre cada usina (linha cheia) e a mesma sob céu claro (tracejado). \textbf{c}, Decomposição da queda do fator de capacidade de Sol do Cerrado entre 2023 e 2025: observado em 2023, esperado em 2025 caso apenas a irradiância tivesse mudado, e observado em 2025. \textbf{d}, Fator de capacidade mensal contra irradiância mensal em Sol do Cerrado, por ano. Dados climáticos: NASA POWER (MERRA-2 + SYN1DEG), série diária no ponto de cada usina. Dados de energia: ONS. Dados-fonte em \texttt{data/processed/fig5\_source\_*.csv}.}
\newcommand{\legVI}{\textbf{Fig. 6 | A geração de referência do ONS é bem calibrada; o viés aparente da Fig. 5 vinha do contrafactual, não do operador.} \textbf{a}, Distribuição do viés mediano por usina, definido como a razão entre o potencial implicado pelo ONS (verificado + cortado) e o potencial previsto pela irradiância local usando a razão de desempenho calibrada nos dias sem restrição da própria usina; tracejado em 1 marca ausência de viés. \textbf{b}, Viés contra a intensidade mediana de corte nos dias restritos. \textbf{c}, Validação da calibração: fator de capacidade contra irradiância nos dias sem restrição de Sol do Cerrado, com a reta ajustada pela razão de desempenho limpa. \textbf{d}, Razão de desempenho anual de Sol do Cerrado contra o desempenho limpo calibrado; o rótulo interno indica o corte implícito de cada ano. Dados: ONS (fator de capacidade, constrained-off) e NASA POWER. Dados-fonte em \texttt{data/processed/fig6\_source\_*.csv}.}
\newcommand{\legVII}{\textbf{Fig. 7 | A construção de usinas fotovoltaicas é detectável e datável por Sentinel-2; a operação não é.} \textbf{a}, Razão entre o brilho no vermelho dentro do arranjo e num anel de controle de 1,5 a 3 km, por usina (linhas cinza) e mediana entre usinas (linha vermelha, faixa entre quartis), alinhadas pela data de entrada em operação do SIGA. \textbf{b}, Sentinel-2 em cor natural sobre o complexo Sol do Cerrado, em escala de refletância fixa entre as datas; círculos marcam as 17 usinas do SIGA. \textbf{c}, Comparação pareada da razão de brilho entre fases, uma linha por usina, com a mediana em destaque; \texttt{p} de Wilcoxon pareado. \textbf{d}, Diferença entre o ano de pico do brilho relativo e o ano de entrada em operação; barras em vermelho indicam erro de até um ano. Imagens: Sentinel-2 L2A (ESA), via catálogo STAC público. Registro: ANEEL/SIGA. Dados-fonte em \texttt{data/processed/fig7\_source\_*.csv}.}
\newcommand{\legVIII}{\textbf{Fig. 8 | O corte fotovoltaico por excesso de oferta coincide com térmica inflexível, mas apenas em carga baixa.} \textbf{a}, Perfil diurno médio nacional. Área azul, geração fotovoltaica verificada; área vermelha empilhada, energia cortada por razão energética (\texttt{ENE}); linhas, geração térmica separada em parcela econômica (por ordem de mérito, quando o CVU da usina é menor que o CMO) e parcela fora do mérito (geração total menos ordem de mérito). A anotação registra a carga líquida --- carga do SIN menos geração fotovoltaica --- na madrugada e no pico solar. \textbf{b}, Composição da geração térmica fora do mérito nas horas com e sem corte por razão energética, decomposta pelos motivos não-econômicos disjuntos do dicionário do ONS; rótulo indica o total. \textbf{c}, Correlação de Spearman entre corte por razão energética e geração térmica fora do mérito, estratificada por quintil de carga do SIN, restrita às horas de 9 h a 15 h; barras cinza indicam associação não significativa; abaixo do eixo, o corte médio de cada quintil. \textbf{d}, Controle negativo: taxa diária de corte por razão energética contra a energia armazenada do SIN, com a mediana por quartil de armazenamento. Dados: ONS, dados abertos --- restrição por \textit{constrained-off} fotovoltaico, geração térmica por motivo de despacho, curva de carga e energia armazenada por subsistema. Dados-fonte em \texttt{data/processed/fig8\_source\_*.csv}.}
\newcommand{\legIX}{\textbf{Fig. 9 | O piso térmico inflexível está onde o corte não está, e a coincidência nacional da Fig. 8 não sustenta leitura causal local.} \textbf{a}, Potência média nas horas solares (9 h--15 h) por subsistema: corte fotovoltaico por razão energética, geração térmica fora da ordem de mérito e a parcela desta que é inflexibilidade pura; abaixo do eixo, a razão entre o corte e o piso inflexível local. \textbf{b}, Correlação de Spearman entre corte por razão energética e térmica fora do mérito \textbf{do mesmo subsistema}, por quintil de carga do próprio subsistema; tracejado no limiar de magnitude declarado. \textbf{c}, A mesma correlação medida contra a térmica do próprio subsistema e contra a do outro subsistema com corte material. \textbf{d}, Distribuição acumulada do fluxo Nordeste$\rightarrow$Sudeste nas horas de corte alto (quartil superior do corte do Nordeste) e nas demais; tracejado no percentil 95 do fluxo, usado como referência empírica de teto por não haver limite operativo publicado. Dados: ONS, dados abertos --- \textit{constrained-off} fotovoltaico, geração térmica por motivo de despacho, curva de carga e intercâmbio nacional. Dados-fonte em \texttt{data/processed/fig9\_source\_*.csv}.}
\newcommand{\legX}{\textbf{Fig. 10 | A distância de rede até a demanda explica parte da alocação do corte que covariáveis locais não explicavam.} \textbf{a}, Grafo de transmissão em operação --- subestações e linhas da EPE, separadas em 500 kV ou mais e abaixo disso --- com as usinas fotovoltaicas da amostra posicionadas no seu ponto de conexão; a cor codifica a taxa de corte na janela de referência da frota e o tamanho, a capacidade instalada. Estrelas marcam as nove regiões metropolitanas usadas como centros de carga; a moldura cobre a região que contém a amostra, e 353 das 1.896 linhas do grafo ficam fora dela e não são desenhadas --- o recorte é de desenho, e todas as arestas participam do cálculo dos caminhos mínimos. \textbf{b}, Correlação de Spearman entre a taxa de corte e cada uma das oito covariáveis de rede pré-declaradas; em vermelho as que sobrevivem à correção de Holm para oito hipóteses. \textbf{c}, Taxa de corte contra a distância de rede ao centro de carga mais próximo, por subsistema, com a mediana por quartil de distância. \textbf{d}, As duas covariáveis sobreviventes medidas \textbf{dentro} de cada subsistema, para separar efeito de rede de efeito de subsistema. Dados: EPE, camadas de linhas de transmissão e subestações em operação; ONS, restrição por \textit{constrained-off} fotovoltaico e fator de capacidade. Dados-fonte em \texttt{data/processed/fig10\_source\_*.csv}.}
\newcommand{\legXI}{\textbf{Fig. 11 | O achado da Fig. 10 não se replica em eólica, e a dissociação espelha a origem do corte de cada fonte.} \textbf{a}, Correlação de Spearman entre a taxa de corte e as oito covariáveis de rede pré-declaradas na Fig. 10, medidas separadamente em fotovoltaica e eólica; a cor cheia marca as que sobrevivem à correção de Holm aplicada dentro de cada fonte, e o cinza as que não sobrevivem. \textbf{b}, Energia cortada decomposta por razão e origem da restrição, conforme classificação operativa do ONS, para cada fonte. \textbf{c}, As duas covariáveis-chave medidas apenas no Nordeste, onde as duas amostras coexistem, para separar tecnologia de geografia. \textbf{d}, Taxa de corte eólica contra o número de conjuntos no mesmo ponto de conexão, com a mediana por valor. Dados: ONS, restrição por \textit{constrained-off} eólico e fotovoltaico e fator de capacidade; EPE, linhas de transmissão e subestações em operação. Dados-fonte em \texttt{data/processed/fig11\_source\_*.csv}.}
\newcommand{\legXII}{\textbf{Fig. 12 | Dos quatro achados de rede da série, um passa nos três testes de robustez e o principal da Fig. 10 passa em apenas um.} \textbf{a}, Distribuição bootstrap do coeficiente de Spearman de cada achado, com 4.000 reamostras com reposição; a barra marca o intervalo de 95\% e o ponto, o valor da amostra. \textbf{b}, Coeficiente calculado separadamente nas duas metades de uma partição estrutural --- usinas cuja barra de conexão consta do cadastro de rede de 230 kV ou mais do ONS, e as demais --- com o valor $p$ do teste $z$ de Fisher para a diferença entre metades. \textbf{c}, O mesmo coeficiente calculado no grafo da EPE, usado nas Figs. 10 e 11, e no grafo do ONS, construído por número de barra; abaixo de cada barra, o tamanho da amostra no grafo alternativo, ou a marca de que o critério não se aplica --- a capacidade da usina não deriva do grafo e não tem versão alternativa. \textbf{d}, Placar: quantos dos critérios aplicáveis cada achado passa; o traço marca o critério não aplicável. Dados: Figs. 10 e 11; ONS, linhas de transmissão e subestações da Rede de Operação e fator de capacidade. Dados-fonte em \texttt{data/processed/fig12\_source\_robustez.csv}.}